\subsubsection{Agmon 条件与强制问题}\label{subsec:agmon-condition-coercive-problems}

考虑 $2m$ 阶标量椭圆算子, 如 \eqref{eq:general-even-order-divergence-form-definition} 那样写为散度形式:
\begin{equation}\label{eq:agmon-coercive-divergence-operator}
L(\bm{x},D)u=\sum_{|\alpha|,|\beta|\leq m}(-1)^{|\alpha|}D^\alpha(a_{\alpha\beta}(\bm{x})D^\beta u),
\end{equation}
其中 $a_{\alpha\beta}$ 在 $\overline\Omega$ 上连续, 且椭圆性条件
\begin{equation}\label{eq:agmon-principal-ellipticity-condition}
\sum_{|\alpha|=|\beta|=m}\bm{\xi}^\alpha a_{\alpha\beta}(\bm{x})\bm{\xi}^\beta>0
\end{equation}
在整个 $\overline\Omega$ 上成立. 此外, 考虑 $p$ 个正规边值算子 $B_j(\bm{x},D)$, 其系数属于 $C^{m-m_j}(\partial\Omega)$ 类, 其中 $m_j<m$ 为 $B_j$ 的阶数. 一般地, $p$ 可以取 $0$ 到 $m$ 之间的任意值. 定义
\begin{equation}\label{eq:agmon-essential-boundary-admissible-space}
V=\{u\in H^m(\Omega)\mid B_j(\bm{x},D)u=0\text{ on }\partial\Omega,\ j=1,\ldots,p\}.
\end{equation}


考虑二次型
\begin{equation}\label{eq:agmon-boundary-constrained-quadratic-form}
a(u,v)=\int_\Omega\sum_{|\alpha|\leq m,|\beta|\leq m}a_{\alpha\beta}(\bm{x})D^\beta uD^\alpha v\,\mathrm{d}\bm{x},
\end{equation}
并询问在什么条件下这个型在 $V$ 上强制:
\begin{equation}\label{eq:agmon-boundary-constrained-coercive-estimate}
a(u,u)\geq c_1\|u\|_m^2-c_2\|u\|_0^2\quad\forall u\in V.
\end{equation}


若该型强制, 则应用 Lax--Milgram 引理可知, 对足够大的 $\lambda$, 方程
\begin{equation}\label{eq:agmon-shifted-constrained-variational-equation}
a(u,v)+\lambda(u,v)=(f,v)\quad\forall v\in V
\end{equation}
对每个 $f\in V\prime$ 都有唯一解 $u\in V$. 于是显然有 $L(\bm{x},D)u+\lambda u=f$ 在分布意义下成立, 且在边界上有 $B_j(\bm{x},D)u=0$. 此外, $u$ 将满足 $m-p$ 个 ``自然\textquotesingle\textquotesingle\ 边界条件, 它们的产生方式与 Section~\ref{subsec:neumann-problem} 中的 Neumann 边界条件类似.

保证强制性的条件称为 Agmon 条件. 考虑点 $\bm{x}_0\in\partial\Omega$; 可以选取坐标系, 使 $\bm{x}_0$ 为原点, 且内法向指向 $x_n$ 方向. 然后考虑半空间 $x_n>0$ 中的常系数问题 $L^p(0,D)u=0$, 并对 $j=1,\ldots,p$ 施加边界条件 $B_j^p(0,D)u=0$. 使用记号 $\bm{x}=(\bm{x}^{\prime},x_n)$, 其中 $\bm{x}^{\prime}\in\mathbb{R}^{n-1}$; 相应地, 对多重指标 $\alpha$ 写成 $\alpha=(\alpha^{\prime},\alpha_n)$. 现在任取 $\bm{\xi}^{\prime}\in\mathbb{R}^{n-1}\backslash\{0\}$, 考虑 ODE
\begin{equation}\label{eq:agmon-principal-half-line-ode}
L^p\left(0,i\bm{\xi}^{\prime},\frac{\mathrm{d}}{\mathrm{d}t}\right)v(t)=0,\ t>0,
\end{equation}
及初始条件
\begin{equation}\label{eq:agmon-principal-half-line-boundary-conditions}
B_j^p\left(0,i\bm{\xi}^{\prime},\frac{\mathrm{d}}{\mathrm{d}t}\right)v(0)=0,\quad j=1,\ldots,p.
\end{equation}


\begin{Definition}
\label{def:agmon-decaying-half-line-positivity}
若对任意 $\bm{\xi}^{\prime}\in\mathbb{R}^{n-1}\backslash\{0\}$, 以及满足 \eqref{eq:agmon-principal-half-line-boundary-conditions} 的 \eqref{eq:agmon-principal-half-line-ode} 的任意非零解 $v$, 只要 $v$ 在 $t\to\infty$ 时指数趋于零, 就有不等式
\begin{equation}\label{eq:agmon-decaying-principal-half-line-positivity}
\int_0^\infty\sum_{\substack{|\alpha^{\prime}|+k=m\\|\beta^{\prime}|+l=m}}a_{(\alpha^{\prime},k)(\beta^{\prime},l)}(0)(\bm{\xi}^{\prime})^{\alpha^{\prime}}(\bm{\xi}^{\prime})^{\beta^{\prime}}\frac{\mathrm{d}^kv(t)}{\mathrm{d}t^k}\frac{\mathrm{d}^lv(t)}{\mathrm{d}t^l}\,\mathrm{d}t>0.
\end{equation}
成立, 则称 Agmon 条件成立.
\end{Definition}

\begin{Remark}
\label{rem:agmon-dirichlet-vacuity}
若 $p=m$ 且补条件成立, 则 Agmon 条件空真. 事实上, 若 $p=m$, 则由 Lemma~\ref{lem:dirichlet-system-tangential-triangular-reduction}, 这些边界条件等价于 Dirichlet 条件. 实际上, Dirichlet 条件总满足补条件; 见 Problem~\ref{ex:scalar-dirichlet-complementing-condition}.
\end{Remark}
